\section{Applications of Laplace}

\subsection{Transfer Functions}

\subsubsection{Transfer function from differential equation}
\[
\left(\sum_{k=0}^{N}a_ks^k\right)Y(s)=
\left(\sum_{m=0}^{M}b_ms^m\right)X(s),
\qquad
H(s)=\frac{Y(s)}{X(s)}=\frac{P(s)}{Q(s)}.
\]
With zero initial conditions, derivatives are replaced by powers of \(s\). The rational transfer function describes the zero-state system.

\subsubsection{Frequency response from transfer function}
\[
H(j\omega)=H(s)\big|_{s=j\omega}.
\]
For stable systems, the Fourier-domain transfer function is obtained by evaluating the Laplace transfer function on the imaginary axis.

\subsubsection{Poles and zeros}
\[
H(s)=K\frac{\prod_i(s-z_i)}{\prod_k(s-p_k)}.
\]
Zeros \(z_i\) make the transfer function vanish. Poles \(p_k\) determine natural modes and stability.

\subsubsection{Partial fraction expansion}
\[
\frac{B(s)}{\prod_{k=1}^{N}(s-p_k)}
=\sum_{k=1}^{N}\frac{A_k}{s-p_k},
\qquad
A_k=\left.(s-p_k)H(s)\right|_{s=p_k}.
\]
For distinct poles, residues give the time-domain exponential weights.

\subsection{Circuit Relations}

\subsubsection{Resistor, capacitor, and inductor laws}
\[
v_R(t)=Ri_R(t),\qquad i_C(t)=C\frac{dv_C(t)}{dt},\qquad v_L(t)=L\frac{di_L(t)}{dt}.
\]
These component equations convert circuits into differential equations in the time domain.

\subsubsection{Laplace impedances}
\[
Z_R=R,\qquad Z_C=\frac{1}{sC},\qquad Z_L=sL.
\]
With zero initial conditions, circuit analysis in the Laplace domain uses algebraic impedances.

\subsubsection{Voltage divider transfer function}
\[
H(s)=\frac{Y(s)}{X(s)}=\frac{Z_2(s)}{Z_1(s)+Z_2(s)}.
\]
Voltage dividers are a direct way to identify first-order RC and RL filter transfer functions.

\subsubsection{First-order RC lowpass and highpass}
\[
H_{\mathrm{LP}}(s)=\frac{1}{1+sRC},\qquad
H_{\mathrm{HP}}(s)=\frac{sRC}{1+sRC}.
\]
The cutoff angular frequency is \(\omega_c=1/(RC)\). The lowpass passes DC, while the highpass rejects DC.

\subsection{Response Calculation}

\subsubsection{Zero-state response in Laplace domain}
\[
Y_{\mathrm{zs}}(s)=H(s)X(s).
\]
The zero-state response is found by multiplying input transform and transfer function, then inverting.

\subsubsection{Zero-input response from initial conditions}
\[
Q(s)Y(s)-\text{initial-condition terms}=0.
\]
The zero-input response is found by setting the input to zero and retaining the unilateral derivative initial-condition terms.

\subsubsection{Step response final value}
\[
y_{\infty}=\lim_{s\to0}sH(s)\frac{1}{s}=H(0).
\]
For a stable system, the final value of the unit step response equals the DC gain.
